We now turn to the general argument for $d\ge3$.  Consider an 
$r$-edge-coloured product of complete graphs.  We say a $1$-dimensional $K_k$ is $1$-\textit{consistent} if it is monochromatic. For $d\geq 2$, we say $\ourbox^{d}\!K_k$ is $d$-\textit{consistent} if the following two conditions hold:
\begin{itemize}
\item for every $a\in [k]$, the subgraph $(\ourbox^{d-1}\!K_k)\Osq a$ has the same colouring: that is, for every edge $\textbf{x}\textbf{y}$ in $\ourbox^{d-1}\!K_k$, the colour of the edge between $\textbf{x}\times a$ and $\textbf{y}\times a$ is the same for every $a\in [k]$; and 
\item for some (and thus for every) $a$, the subgraph $\ourbox^{d-1}\!K_k\Osq a$ is $(d-1)$-\textit{consistent}. 
\end{itemize}
In other words, for each $i\in[d]$, the `$i$-dimensional subspaces' $(\ourbox^i K_N)\Osq a_{i+1}\Osq\cdots\Osq a_d$ have the same colouring for every choice of $a_{i+1},\dots,a_d$.

The proof splits into two lemmas. Given suitable $N\gg k\gg n$ and an $r$-edge colouring of $\ourbox^d\!K_N$, we first find a  $d$-\textit{consistent}  subgraph $H=\ourbox^{d}\!K_k$. We then show that $H$ contains a monochromatic copy of $\ourbox^{d}\!K_n$. 


\begin{lemma}\label{lem: 1}
For every $d\ge1$ there is a constant $g(d)$ such that the following holds.  Let $r,k$ be positive integers and $0<\epsilon<1/2$. Suppose that 
$N\geq \epsilon^{-g(d)k^{d-1}}\cdot r^{g(d)rk^{d}}$ and $\ourbox^d K_N$ has a $d$-consistent $r$-edge-colouring. Then every set $S$ of at least $\epsilon N^{d}$ vertices contains a copy of $\ourbox^{d}\!K_k$ which is monochromatic in every direction. 
\end{lemma}

\begin{proof}
We argue by induction on $d$ that $g(d)=(12+2r)^{d-1}$ will do. It is clear that when $d=1$ it is sufficient to have $N\geq r^{rk}/\epsilon$, which we do as $g(1)=1$. So we assume that $d\geq 2$ and we have handled smaller cases. Let $T$ be the set of elements ${\mathbf v}\in[N]^{d-1}$ such that ${\mathbf v}\times[N]$ contains at least $\epsilon N/2$ elements of $S$. A counting argument shows that $|T|\ge \epsilon N^{d-1}/2$.
Let $A\subset[N]$ be a random subset of size $(10/\epsilon)r^{rk}$. For each ${\mathbf v}\in T$, with probability at least $2/3$ the set $\mathbf v\times A$ contains at least $r^{rk}$ elements of $S$. So we can choose $A$ so that the set
    $$T':=\{ {\mathbf v}\in T: |({\mathbf v}\times A)\cap S|\ge r^{rk}\}$$
has size at least $(2/3)|T|\ge(\epsilon/3)N^{d-1}$.
     By Ramsey's Theorem, for each ${\mathbf v}\in T'$, there is a monochromatic copy of $K_k$ contained in 
$({\mathbf v}\times A)\cap S$, 
say on vertices $B_{\mathbf v}\subset A$. There are at most $r\binom{(10/\epsilon)r^{rk}}{k}\leq \epsilon^{-9k}r^{rk^2+1}$ 
     choices for $B_{\mathbf v}$ and the colour of the corresponding copy of $K_k$, so there are $B\subset A$ and $U\subset T'$ such that
$$|U|\ge \epsilon^{9k}r^{-rk^2-1}|T'| \ge (\epsilon/3)\epsilon^{9k}r^{-rk^2-1}N^{d-1}$$
and $B_{\mathbf v}=B$ for every $\mathbf v\in U$.  
Now let $\tilde{\epsilon}=(\epsilon/3)\epsilon^{9k}r^{-rk^2-1}\ge \epsilon^{12k}r^{-rk^2-1}$, so $|U|\geq \tilde{\epsilon}N^{d-1}$. We now apply the inductive hypothesis to $[N]^{d-1}$, with the colouring inherited from $[N]^{d-1}\times a$ (which by consistency is the same for any $b\in B$) with the subset $U$ playing the role of $S$.
Since $N\geq \epsilon^{-g(d)k^{d-1}}r^{g(d)rk^d}
\geq \tilde{\epsilon}^{-g(d-1)k^{d-2}}\cdot r^{g(d-1)rk^{d-1}}$, 
we obtain a $(d-1)$-dimensional product $P$ of copies of $K_k$ which is monochromatic in every direction.  Then $P\times B$ gives a $d$-dimensional product of copies of $K_k$ which is monochromatic in every direction (as $P\times b$ is coloured in the same way for every $b\in B$, and all sets $\textbf{v}\times B$ ($\textbf{v}\in B)$ give monochromatic copies of $K_k$ with the same colour).
\end{proof}


\begin{lemma}\label{lem: 2}
For every positive integer $d$ there is $f(d)$ such that for every $r,k$ the following holds. Let $N\coloneqq r^{f(d)rk^{d}}$. Then in every $r$-colouring of $\ourbox^{d}\!K_N$ there is a $d$-consistent $\ourbox^{d}\!K_k$. 
\end{lemma}
\begin{proof}
We argue by induction on $d$ that $f(d)=2^{d-1}d!$ will do. For $d=1$, this is true by Ramsey's Theorem. 
Now, let $M\coloneqq r^{f(d-1)rk^{d-1}}$ and consider the subgraph $(\ourbox^{d-1}\!K_{M})\Osq K_N$. By induction any $r$-colouring of $\ourbox^{d-1}\!K_{M}$ contains a $(d-1)$-consistent $\ourbox^{d-1}\!K_k$. Therefore, for every $a\in [N]$ the subgraph $\ourbox^{d-1}\!K_{M}\times a$ contains a $(d-1)$-consistent $\ourbox^{d-1}\!K_k\times a $.
 As there are most $r^{(d-1)k^{2+(d-2)}}$ possible $r$-colourings of $\ourbox^{d-1}K_k$ and at most $\binom{M}{k}^{d-1}$ possible vertex sets, 
 there are at most $r^{(d-1)k^{2+(d-2)}}\binom{M}{k}^{d-1}\le r^{(d-1)k^d}M^{k(d-1)}/k< r^{f(d)rk^d}/k=N/k$ possible combinations.
Thus, by the pigeonhole principle, there is a $(d-1)$-consistent colour pattern $c$ of $\ourbox^{d-1}K_k$ and a set $A\subset [N]$ of size $k$, such that for every $a\in A$, the subgraph $\ourbox^{d-1}\!K_k \times a$ has colour pattern $c$, and all these boxes lie on the same vertex set in the first $d-1$ dimensions, as we wanted to show. 
\end{proof}

\begin{proof}[Proof of Theorem~\ref{thm:mainRamsey}]
We will show that the statement holds with $C_d=3dg(d-1)+f(d)+1$, where $f$ and $g$ are
any functions satisfying the previous two lemmas.
Let $N\coloneqq r^{r^{C_drn^{d}}}$. Let $t\coloneqq r^{3g(d-1)rn^{d}}$ and $u\coloneqq r^{rn}$. 
Applying Lemma~\ref{lem: 2} to $\ourbox^{d}\!K_{N}$, we obtain a $d$-consistent copy $B$ of $\ourbox^{d}\!K_t$. Relabelling, and restricting the final coordinate to $u$ choices, we may assume that we have a $d$-consistent colouring of $\ourbox^{d-1}\!K_t\Osq K_u$. 

For every $a_1\times \dots \times a_{d-1}$, where each $a_i\in [t]$, we can apply Ramsey's Theorem to $a_1\Osq \dots \Osq a_{d-1}\Osq K_{u}$ to get a monochromatic copy of $K_n$. 
There are $\binom{u}{n}\cdot r$ choices of colour and coordinates $C$ for this copy; setting
$\epsilon=1/(ru^{n})$ we see that there is some set $S$ of at least $\epsilon t^{d-1}$ vertices in $\ourbox^{d-1}\!K_t$ for which both colour and coordinates agree.

%
%
%

Finally, we apply \Cref{lem: 1} to 
$S\subset\ourbox^{d-1}\!K_t$
and obtain a copy $K$ of $\ourbox^{d-1}\!K_n$ contained in $S$ that is monochromatic in every direction. This lemma holds for
$t\ge \epsilon^{-g(d-1)n^{d-2}}r^{g(d-1)rn^{d-1}}$, which holds by our choice of constants,
as 
$$\epsilon^{-g(d-1)n^{d-2}}
=(ru^n)^{g(d-1)n^{d-2}}
=r^{(rn^2+1)g(d-1)n^{d-2}}.$$
Then $K\Osq C$ is a copy of $\ourbox^{d}\!K_n$ that is monochromatic in every direction.
%
\end{proof}
